As Large Language Models (LLMs) continue to advance, efficient inference remains a challenge. Autoregressive decoding generates one token at a time. This inherent seriality creates a major performance bottleneck: decoding is slow, memory-bound, and fails to fully utilize modern GPUs~\citep{cai2024medusa}.

Speculative decoding accelerates a language model, called the \emph{target}, using a smaller model or prediction module, called the \emph{drafter}. The drafter proposes several candidate tokens, and the target checks them in parallel. Verification keeps an accepted prefix and replaces the first rejected proposal with a correction. Appropriate acceptance and correction rules preserve the target's output distribution~\citep{leviathan2023speculative,chen2023sampling}. Speedup depends primarily on the length of the accepted prefix in each proposed block, together with the time spent drafting and verifying that block.

Some drafters also use the target's hidden-state representations. These representations provide contextual information that can help the drafter propose tokens the target will accept. Medusa adds parallel prediction heads to the target's final hidden states~\citep{cai2024medusa}. The EAGLE family uses target features to guide autoregressive drafting~\citep{li2024eagle,li2024eagle2,li2025eagle3}. DFlash instead uses a block-diffusion drafter to predict several tokens in one pass and injects fused target hidden-state representations into every layer of the drafter model~\citep{chen2026dflash}.

This dependence complicates reuse of the drafter model when a service changes its target's size or family, or LoRA-fine-tunes it~\citep{hu2022lora}. The new target can produce different hidden states and preferred continuations, and changed feature dimensions may be incompatible with the existing drafter. Training a new drafter model requires target-specific supervision, optimization of the drafter's weights, and validation~\citep{li2025eagle3,chen2026dflash}. Repeating this process across target variants increases training and maintenance costs. Can we instead reuse an existing drafter by adapting only its target-feature inputs?

We introduce RelaySpec, a method for reusing a pretrained DFlash drafter~\citep{chen2026dflash} after its target changes. Section~\ref{sec:method} explains the parts of DFlash that RelaySpec uses. The new target supplies hidden states from five layers. A separate linear map transforms each stream before the drafter's original fusion projection combines them. The resulting context conditions the drafter transformer, which predicts a block of candidate tokens.

During adaptation, only the five maps are trained. The target, original fusion projection, drafter transformer, normalization layers, and the embeddings and output head selected for the model pair remain frozen. A map can be rectangular when the target and drafter expect different feature dimensions. Before inference, we multiply each map into its corresponding fusion block, a composition of linear operations. This produces one effective projection, so inference needs no separate mapping operation.

RelaySpec uses training signals that match the change in the feature interface. For a LoRA-fine-tuned target, identity-initialized maps preserve the existing interface and are trained to improve the accepted prefix. For cross-model transfer, the maps reconstruct the original target's features and fused context from the new target's features. Section~\ref{sec:method} gives both procedures.

Figures~\ref{fig:lora-ar} and~\ref{fig:transfer-ar} show the measured throughput in these two settings relative to autoregressive decoding.

We analyze the learned maps by compressing their corrections with SVD, scaling the correction strength, comparing the corrections with LoRA-induced feature changes, and interpolating between token-trained and MSE-trained maps. We measure how these interventions change acceptance and throughput.
